%%%%%%%%%%%%%%%%%%%%%%%%%%%%%%%%%%%%%%%%%%%%%%%%%%%%%%%%%%%%%%%%%%%%%%
% Lớp 11 | Đường thẳng, mặt phẳng. Quan hệ song song trong | Hình lăng trụ và hình hộp (xiên)
% Dạng 1H4?5-1: Câu hỏi lý thuyết
% Số câu: 20  (sinh tự động bởi tach_cau.py; từ nay thầy sửa trực tiếp file này)
%%%%%%%%%%%%%%%%%%%%%%%%%%%%%%%%%%%%%%%%%%%%%%%%%%%%%%%%%%%%%%%%%%%%%%

%%%%%%%%%%%%%%%% TRẮC NGHIỆM 4 PHƯƠNG ÁN %%%%%%%%%%%%%%%%

% nguồn: Ok_Toan11_TracNghiem_Lan2
\begin{ex}%[1H4N5-1]%[ma:61e30e2e]
	Trong các khẳng định sau, khẳng định  nào \textbf{sai}?
	 \choice
		 {Hai đáy của hình lăng trụ là hai đa giác bằng nhau}
		{\True Các mặt bên của hình lăng trụ là các  hình bình hành bằng nhau}
		{Các cạnh bên của hình lăng trụ bằng  nhau và song song với nhau}
		{Các mặt bên  của hình lăng trụ là các hình bình hành}
	 \loigiai{
		\begin{center}
			\begin{tikzpicture}[scale=0.8,font=\footnotesize,line join=round,line cap=round,>=stealth]
					\coordinate (A) at  (0,0);
					\coordinate (B) at  (2,-1);
					\coordinate (C) at (2,0); 
					\coordinate  (D) at (3,2);
					 \coordinate (E) at (1,2);
					\coordinate (C) at  ($(B)+(D)-(A)$);
					\coordinate (H)  at (1,4);
					\foreach \i in {A,B,C,D,E}{\coordinate (\i') at  ($(\i)+(H)$);}
					 \draw (A)--(B)--(C);
					 \draw (A')--(B')--(C')--(D')--(E')--cycle;
					\draw  (A)--(A')
						(B)--(B') 
						 (C)--(C');
					\draw[dashed,thin](A)--(E)--(D)--(C) 
						 (D)--(D')
						 (E)--(E');
					\foreach \i/\g in  {A'/90,B'/90,C'/90,D'/90,A/-90,B/-90,C/-90,D/-90,E/90,E'/90}{\draw[fill=blue](\i) circle (1pt)
						 ($(\i)+(\g:5mm)$) node[scale=1]{$ $};}
				\end{tikzpicture}
		\end{center}
		Theo tính chất của hình  lăng trụ thì\\
		Hai đáy của hình lăng trụ là hai đa  giác bằng nhau và nằm trên hai mặt phẳng song song nhau.\\
		Các cạnh bên  của hình lăng trụ bằng nhau và song song với nhau.\\
		Các mặt bên của hình lăng trụ là các hình bình hành nhưng chưa chắc  chúng bằng nhau.
	} 
\end{ex}

% nguồn: Ok_Toan11_TracNghiem_Lan2
\begin{ex}%[1H4N5-1]%[ma:ab283855]
	Cho hình lăng trụ tam  giác, chọn khẳng định \textbf{sai}.
	 \choice
		{Các cạnh bên song song  với nhau}
		{\True Các  mặt bên là các hình chữ nhật}
		{Hai tam  giác đáy bằng nhau}
		{Hai đáy nằm trên hai mặt  phẳng song song}
	\loigiai{ 
		\begin{center}
			\begin{tikzpicture}[scale=0.8,font=\footnotesize,line join=round,line cap=round,>=stealth]
					\coordinate (A)  at (-2,0);
					\coordinate (B) at (-0.5,-2); 
					\coordinate (C) at (2,0); 
					\foreach \i in {A,B,C}{\coordinate  (\i') at ($(\i)+(1,4)$); \draw (\i)--(\i');}
					\draw  (A')--(B')--(C')--cycle;
					 \draw (A)--(B)--(C);
					 \draw[dashed,thin](A)--(C);
					\foreach \i/\g in  {A/180,B/-90,C/0,A'/180,B'/-30,C'/0}{\draw[fill=black](\i) circle (1pt)
						($(\i)+(\g:4mm)$)  node[scale=1]{$\i$};}
				\end{tikzpicture}
		\end{center}
		Theo định nghĩa hình lăng trụ thì  các mặt bên là hình bình hành.
		Vậy khẳng định \lq \lq  Các mặt bên là hình chữ nhật\rq\rq là sai.
	} 
\end{ex}

% nguồn: Ok_TracNghiemLop11_Lan1
\begin{ex}%[1H4N5-1]%[ma:c8ac5daf]
	Hình lăng trụ có các mặt bên là hình gì? 
	 \choice
		{Hình thoi} 
		{Hình  vuông}
		{Hình chữ  nhật}
		{\True Hình bình hành} 
	\loigiai{Hình lăng trụ có các  mặt bên là các hình bình hành.
	 }
\end{ex}

% nguồn: Ok_TracNghiemLop11_Lan1
\begin{ex}%[1H4N5-1]%[ma:59ad8194]
	Cho hình hộp $ABCD.A'B'C'D'$. Trong các khẳng định sau, khẳng  định nào \textbf{sai}?
	\choice 
		 {$\left(ABCD\right)\parallel \left(A'B'C'D'\right)$}
		 {$\left(AA'D'\right)\parallel \left(BCC'\right)$}
		{\True  $\left(BDD'\right)\parallel \left(ACC'\right)$}
		{$\left(ABB'\right)\parallel \left(CDC'\right)$} 
	 \loigiai{
	\immini{Gọi $O$ là giao điểm của $AC$ và $BD$.  \\
		Khi đó, $O\in (BDD')\cap (ACC')$ nên $(BDD')$ không song song với $(ACC')$. }{\begin{tikzpicture}[scale=1,font=\footnotesize,line join=round,line cap=round,>=stealth]
		\path (-0.5,2) coordinate (h) 
			(0,0)  coordinate (B)++(2.5,0) coordinate (C)++(1,0.8) coordinate (D)
			($(B)+(D)-(C)$) coordinate (A) 
			(barycentric cs:A=1,C=1) coordinate (O); 
		\foreach  \p in {A,B,C,D}\path ($(\p)+(h)$) coordinate (\p');
		\draw (A')--(B')--(B)--(C)--(D)--(D')--cycle (B')--(C')--(D') 
			 (C)--(C');
		 \draw[dashed] (B)--(A)--(D)--(B)
			(C)--(A)--(A'); 
		\foreach  \p/\g in {A/150,B/240,C/-60,D/0,A'/90,B'/180,C'/-30,D'/45,O/-90} \fill[black] (\p) circle (1pt)+(\g:0.3) node{$\p$};
	\end{tikzpicture}
		 }
	 }
\end{ex}

% nguồn: Ok_TracNghiemLop11_Lan1
\begin{ex}%[1H4N5-1]%[ma:e524d302]
	Cho hình lăng  trụ tam giác, chọn khẳng định \textbf{sai}?
	 \choice
		{Các cạnh bên song  song với nhau}
		{\True Các mặt bên  là các hình chữ nhật}
		{Hai tam giác  đáy bằng nhau}
		{Hai đáy nằm trên hai  mặt phẳng song song}
	\loigiai{ 
		Theo định nghĩa hình lăng trụ thì  các mặt bên là hình bình hành. Vậy khẳng định \lq \lq Các mặt bên là hình chữ nhật\rq \rq \, là \textbf{sai}.
	 }
\end{ex}

% nguồn: Ok_TracNghiemLop11_Lan1
\begin{ex}%[1H4N5-1]%[ma:b8c996e4]
	Lăng trụ ngũ  giác có số đỉnh, số cạnh, số mặt lần lượt là $n$, $ p$, $ q$. Chọn khẳng định đúng
	 \choice
		{$n=5$;  $p=15$; $q=7$}
		{$n=5$;  $p=10$; $q=7$}
		{$n=10$; $p=15$; $q=5$} 
		{\True $n=10$; $p=15$; $q=7$} 
	 \loigiai{
		\begin{itemize}
			\item Số đỉnh $n=5\cdot  2=10$.
			 \item Số cạnh $p=5\cdot 2 +5=15$.
			\item  Số mặt $q=2+5=7$.
		\end{itemize}
		Vậy, lăng trụ ngũ giác có số đỉnh, số cạnh, số mặt lần lượt là $n=10$, $p=15$  và $q=7$.
	} 
\end{ex}

% nguồn: Ok_TracNghiemLop11_Lan1
\begin{ex}%[1H4N5-1]%[ma:0f174306]
	Tìm khẳng  định \textbf{ sai } về hình lăng trụ.
	\choice 
		{ \True Hình lăng trụ tứ giác  có mặt bên là hình tứ giác}
		{ Hình lăng trụ có các cạnh bên song song} 
		{ Hình hộp là hình lăng trụ có đáy là  hình bình hành}
		{ Hình lăng  trụ tam giác có đáy là tam giác}
	 \loigiai{
		 \lq\lq Hình lăng trụ tứ giác có mặt bên là hình tứ giác\rq\rq\, là khẳng định sai.
	} 
\end{ex}

% nguồn: Ok_TracNghiemLop11_Lan1
\begin{ex}%[1H4N5-1]%[ma:4662d248]
	Phát  biểu nào dưới đây là đúng về phép chiếu song song?
	 \choice
		{\True Phép chiếu  song song biến hai đường thẳng song song thành hai đường thẳng song song hoặc trùng nhau}
		{Phép chiếu song  song biến hai đường thẳng trùng nhau thành hai đường thẳng song song}
		{Phép chiếu song song biến hai đường thẳng bất kỳ thành hai  đường thẳng song song}
		{Phép chiếu  song song biến hai đường thẳng song song thành hai đường thẳng bất kỳ}
	\loigiai{ 
		Phép chiếu song song biến  hai đường thẳng song song thành hai đường thẳng song song hoặc trùng nhau.
	} 
\end{ex}

% nguồn: Ok_TracNghiemLop11_Lan1
\begin{ex}%[1H4N5-1]%[ma:1d7b1dfb]
	Cho hình hộp $ABCD.A'B'C'D'$. Mệnh đề nào  sau đây \textbf{sai}?
	\choice 
		{$\left(BA'C'\right)\parallel\left(ACD'\right)$} 
		{$\left(ADD'A'\right)\parallel\left(BCC'B'\right)$} 
		 {$\left(BA'D\right)\parallel\left(CB'D'\right)$}
		 {\True $(ABA')\parallel(CB'D')$}
	\loigiai{ 
		\begin{center}
			\begin{tikzpicture}[scale=0.6,font=\footnotesize,line join=round,line cap=round,>=stealth]
					 \def\bc{4}
					\def\ba{2} 
					 \def\h{4}
					\def\gocB{35} 
					\coordinate[label=below  left:$B$] (B) at (0,0);
					\coordinate[label=above left:$A$] (A)  at (\gocB:\ba);
					\coordinate[label=below:$C$] (C) at (\bc,0); 
					\coordinate[label=right:$D$] (D)  at ($(C)-(B)+(A)$);
					\coordinate (H)  at ($(B)!0.4!(D)$);
					 \coordinate[label=above left:$A'$] (A') at ($(H)+(90:\h)$);
					 \coordinate[label=left:$B'$] (B') at ($(B)-(A)+(A')$);
					\coordinate[label=below  right:$C'$] (C') at ($(C)-(A)+(A')$);
					\coordinate[label=right:$D'$] (D') at  ($(D)-(A)+(A')$);
					\draw (B')--(B)--(C)--(D)--(D')--(A')--(B')--(C')--(D') 
						 (C)--(C');
					 \draw[dashed] (A)--(D)
						 (A')--(A)--(B);
					 \foreach \diem in {A,B,C,D,A',B',C',D'} \fill (\diem) circle (1pt);
				\end{tikzpicture}
		\end{center}
		Ta  có\\
		\begin{itemize}
			\item $\heva{&  BA'\parallel CD' \\ & A'C'\parallel AC }$$\Rightarrow (BA'C')\parallel(ACD')$.
			\item $\heva{& AD\parallel BC  \\ & AA'\parallel BB' }$$\Rightarrow (ADD'A')\parallel(BCC'B')$.
			\item $\heva{& BD\parallel B'D' \\ & A'D\parallel B'C }$$\Rightarrow  (BA'D)\parallel(CB'D')$.
		\end{itemize}
	Mặt khác $B'\in (ABA')\cap (CB'D')  $ nên D sai.}
\end{ex}

% nguồn: Ok_TracNghiemLop11_Lan1
\begin{ex}%[1H4N5-1]%[ma:e295a98e]
	Trong các hình sau hình nào biểu diễn cho một hình hộp có  đáy là hình bình hành?
	\begin{center}
		\begin{tikzpicture}[scale=0.6,font=\footnotesize,line join=round,line cap=round,>=stealth]
				 \coordinate (B) at (-2,0);
				\coordinate  (C) at (1,-2);
				 \coordinate (D) at (3,0);
				 \coordinate (A) at ($(B)+(2,3)$);
				 \draw (A)--(B)
					(A)--(C) 
					 (A)--(D)
					(C)--(D) 
					(B)--(C); 
				\draw[dashed](B)--(D); 
				\node at (1,-3) {\textbf{Hình  a}};
			\end{tikzpicture}
		\hspace{1cm} 
		\begin{tikzpicture}[scale=0.4,font=\footnotesize,line join=round,line cap=round,>=stealth]
				\coordinate  (A) at (0,0);
				 \coordinate (B) at (-2,-2);
				\coordinate  (D) at (5,0);
				\coordinate (C) at ($(B)+(D)-(A)$); 
				\coordinate (H) at  (-1,4);
				\foreach \i in {A,B,C,D}{\coordinate  (\i') at ($(\i)+(H)$);}
				 \draw (A')--(B')--(C')--(D')--cycle;
				 \draw (B)--(B')
					(C)--(C') 
					(D)--(D') 
					 (B)--(C)--(D);
				\draw[dashed](B)--(A)--(A') 
					 (A)--(D);
				\node at (1,-3.5) {\textbf{Hình b}}; 
			\end{tikzpicture}
	\end{center}
	 \choice
		 {\True Hình b}
		{Hình a} 
		{Cả  hình a và b}
		{Không có hình nào} 
	\loigiai{ 
	Ta có hình b là hình biểu diễn cho hình hộp  có đáy là hình bình hành}
\end{ex}

% nguồn: Ok_TracNghiemLop11_Lan1
\begin{ex}%[1H4N5-1]%[ma:68430879]
	Cho hình hộp $ABCD.A'B'C'D'$. Cặp điểm  nào sau đây là hai đỉnh đối diện của hình hộp?
	 \choice
		{  $A'$ và $C'$}
		{ \True $A$ và $C'$} 
		{ $A$ và $B'$} 
		{ $A$ và $C$} 
	 \loigiai{
		\begin{center}
			\begin{tikzpicture}[scale=.7,font=\footnotesize,line join=round,line cap=round,>=stealth]
					 \def\bc{4}
					 \def\ba{2}
					\def\gocB{35} 
					\path (0,0) coordinate  (A')
						 (\gocB:\ba) coordinate (B')
						(\bc,0)  coordinate (D')
						($(D')-(A')+(B')$) coordinate  (C')
						($(B')+(80:\bc-1)$)  coordinate (B)
						 ($(A')-(B')+(B)$) coordinate (A)
						($(D')-(B')+(B)$) coordinate  (D)
						($(C')-(B')+(B)$) coordinate  (C);
					\draw (A)--(A')--(D')--(C')--(C)--(B)--(A)--(D)--(C) 
						 (D')--(D);
					\draw[dashed] (B)--(B')--(C') 
						 (B')--(A');
					\foreach \x/\g in {A/90,B/90,C/90,D/90,A'/180,B'/150,C'/0,D'/-30}\fill[black] (\x) circle (1pt)+(\g:0.3)  node{$\x$};
				\end{tikzpicture}
		\end{center}
		Hai  đỉnh đối diện của hình hộp là $A$ và $C'$.
	} 
\end{ex}

% nguồn: Ok_Toan11_TracNghiem_Lan2
\begin{ex}%[1H4H5-1]%[ma:fdc4909a]
	Đặc điểm nào sau  đây là đúng với hình lăng trụ?
	\choice 
		{Hình lăng trụ có tất cả các mặt bên bằng nhau} 
		{Đáy của hình lăng trụ là hình  bình hành}
		 {\True Hình lăng trụ có tất cả các mặt bên là hình bình hành}
		{Hình lăng  trụ có tất cả các mặt là hình bình hành}
	 \loigiai{
	Hình lăng trụ là một đa diện có hai mặt đáy là các đa giác bằng nhau và những mặt bên là các hình bình  hành.}
\end{ex}

% nguồn: Ok_Toan11_TracNghiem_Lan2
\begin{ex}%[1H4H5-1]%[ma:cc2207e4]
	Lăng trụ tam giác  có bao nhiêu mặt?
	\choice 
		 {$3$}
		 {$9$}
		{\True $5$} 
		 {$6$}
	 \loigiai{
		Theo lý thuyết ta có kết quả là  $5$ mặt.
	} 
\end{ex}

% nguồn: Ok_TracNghiemLop11_Lan1
\begin{ex}%[1H4H5-1]%[ma:ecd30165]
	Cho hình lăng trụ $ABC.A'B'C'$. Cắt hình lăng trụ bởi một mặt phẳng ta được một thiết diện. Thiết diện này có tối  đa bao nhiêu cạnh?
	\choice 
		{\True $5$} 
		 {$4$}
		{$3$} 
		{$6$} 
	\loigiai{ 
		Hình lăng trụ $ABC.A'B'C'$ có tất cả $5$ mặt. Do  đó một phẳng cắt hình lăng trụ $ABC.A'B'C'$ theo một thiết diện là hình đa giác có nhiều nhất là $5$ cạnh.
	} 
\end{ex}


%%%%%%%%%%%%%%%% ĐÚNG/SAI %%%%%%%%%%%%%%%%

% nguồn: Ok_dungsai_lop11_Lan2
\begin{ex}%[1H4N5-1]%[ma:cd3bdea8]
	Cho hai hình bình hành $ABCD$ và $ABEF$ nằm trong hai mặt  phẳng phân biệt.
	\choiceTF 
		{$AD\parallel(BEF)$} 
		{\True $(AFD)\parallel(BEC)$} 
		 {$(ABD)\parallel(EFC)$}
		 {\True Sáu điểm $A$, $B$, $C$, $D$, $E$, $F$ là $6$ đỉnh của một hình lăng trụ tam giác}
	\loigiai{ 
		\begin{center}
			\begin{tikzpicture}[line join=round,line cap=round,>=stealth,font=\footnotesize,scale=1]
					 \path
						(0,0)  coordinate (A)
						(0:5) coordinate  (B)
						(-60:2) coordinate  (D)
						(70:2) coordinate (F) 
						 ($(B)+(D)-(A)$) coordinate (C)
						 ($(B)+(F)-(A)$) coordinate (E);
					\draw (A)--(D)--(F)--cycle 
						(D)--(C)--(E)--(F); 
					 \draw[dashed] (A)--(B)--(C)
						 (B)--(E);
					\foreach \i/\g in {A/180,B/0,C/-90,D/-90,E/90,F/90}{\draw[fill=black](\i)  circle (1pt)
						 ($(\i)+(\g:3mm)$) node[scale=1]{$\i$};}
				\end{tikzpicture}
		\end{center}
		\begin{itemchoice}
			\itemch \textbf{Sai}.\\ 
			$AD$ và $(BEF)$ cắt  nhau tại $A$.
			\itemch \textbf{Đúng}.\\ 
			Vì $\heva{&AD\parallel BC\subset (BCE)\\&AF\parallel  BE\subset (BCE)\\&AD\cap AF=A\\&AD,AF \subset (BCE)} \Rightarrow (AFD)\parallel(BEC)$.
			\itemch  \textbf{Sai}.\\
			Vì $(ABD)$  và $(EFC)$ có $C$ chung.
			\itemch \textbf{Đúng}.\\ 
			Vì  $(AFD)\parallel(BEC)$, hai mặt đáy cùng là tam giác; các mặt bên là các hình bình hành.
		\end{itemchoice}
	 }
\end{ex}


%%%%%%%%%%%%%%%% TRẢ LỜI NGẮN %%%%%%%%%%%%%%%%

% nguồn: Ok_TraLoiNgan_Lop11_Lan1
\begin{ex}%[1H4N5-1]%[ma:7897151a]
	Cho hình hộp  $ABCD.A'B'C'D'$. Trong các đoạn thẳng là cạnh hoặc đường chéo của hình hộp, số đoạn thẳng đồng quy tại một điểm là?
	 \shortans{4}
	\loigiai{ 
		\begin{center}
			\begin{tikzpicture}[scale=1,font=\footnotesize,line join=round,line cap=round,>=stealth]
					 \def\bc{4}
					 \def\ba{2}
					\def\h{4} 
					\def\gocB{55} 
					\coordinate[label=below left:$B$] (B) at  (0,0);
					 \coordinate[label=above left:$A$] (A) at (\gocB:\ba);
					\coordinate[label=below:$C$] (C)  at (\bc,0);
					\coordinate[label=right:$D$]  (D) at ($(C)-(B)+(A)$);
					\coordinate (H) at  ($(B)!0.4!(D)$);
					\coordinate[label=above left:$A'$] (A') at  ($(H)+(90:\h)$);
					\coordinate[label=left:$B'$]  (B') at ($(B)-(A)+(A')$);
					 \coordinate[label=below right:$C'$] (C') at ($(C)-(A)+(A')$);
					 \coordinate[label=right:$D'$] (D') at ($(D)-(A)+(A')$);
					\coordinate[label=below:$O$] (O) at  ($(A)!0.5!(C')$);
					\draw (B')--(B)--(C)--(D)--(D')--(A')--(B')--(C')--(D') 
						(C)--(C'); 
					\draw[dashed] (A)--(D) 
						 (A')--(A)--(B)
						 (A')--(C)
						 (B')--(D)
						 (A)--(C')
						 (B)--(D');
					\foreach \diem in {A,B,C,D,A',B',C',D',O} \fill (\diem)  circle (1pt);
				\end{tikzpicture}
		\end{center}
		Dễ thấy $\heva{&AA'\parallel CC'\\ &AA'=CC'}\Rightarrow ACC'A'$ là hình bình hành, suy ra giao điểm $O$ của $AC'$ và $CA'$  là trung điểm của mỗi đường.\\
		Tương tự, ta có $\heva{&AB\parallel C'D'\\ &AB=C'D'}\Rightarrow ABC'D'$ là hình  bình hành; suy ra $ AC'$, $BD'$ cắt nhau tại trung điểm của mỗi đường, do đó $ BD'$ nhận $O$ làm trung điểm. Chứng minh tương tự như thế, ta cũng có $ DB'$ nhận $O$ làm trung điểm.\\
	Vậy các đường chéo của hình hộp là $AC'$, $BD'$, $CA'$, $DB'$ cùng đi qua trung điểm  $O$ của mỗi đường.}
\end{ex}

% nguồn: Ok_TraLoiNgan_Lop11_Lan1
\begin{ex}%[1H4N5-1]%[ma:d7d7e506]
	Trong hình hộp $ABCD.A'B'C'D'$, các đường thẳng $AC'$,  $BD'$, $CA'$ cùng đi qua điểm $I$. Tính tỉ số $\dfrac{IB'}{ID}$.
	 \shortans{1}
	 \loigiai{
	\immini{Trong hình hộp, các đường chéo $AC'$, $BD'$, $CA'$ và $DB'$ đồng  quy tại tâm $I$ là trung điểm của mỗi đường.\\
		Vậy $\dfrac{IB'}{ID}=1$.}{\begin{tikzpicture}[scale=1,font=\footnotesize,line join=round,line cap=round,>=stealth]
		\path (-0.7,2.5) coordinate (h) 
			(0,0) coordinate (B)++(2.2,0) coordinate (C)++(1,0.8) coordinate (D) 
			($(B)+(D)-(C)$)  coordinate (A);
		\foreach \p in {A,B,C,D}\path  ($(\p)+(h)$) coordinate (\p');
		\path (barycentric cs:A=1,C'=1)  coordinate (I);
		\draw  (A')--(B')--(B)--(C)--(D)--(D')--cycle (B')--(C')--(D')
			(C)--(C'); 
		 \draw[dashed] (D')--(B)--(A)--(D)--(B')
			 (C')--(A)--(A')--(C);
		\foreach \p/\g in {A/45,B/240,C/-60,D/0,A'/90,B'/180,C'/90,D'/45,I/10}  \fill[black] (\p) circle (1pt)+(\g:0.3) node{$\p$};
	\end{tikzpicture}
		} 
	} 
\end{ex}


%%%%%%%%%%%%%%%% TỰ LUẬN %%%%%%%%%%%%%%%%

% nguồn: Ok_TuanLuan_Toan11_Lan2
\begin{ex}%[1H4N5-1]%[ma:ff40d2c3]
	\begin{enumerate}
		\item Gọi tên các hình lăng trụ trong các hình sau  đây.
		\item Gọi tên các thành phần của hình lăng trụ trong hình  đầu tiên.
	\end{enumerate}
	\begin{center}
		\begin{tikzpicture}[line join=round,line cap=round,line width=.6pt,font=\footnotesize,scale=.7]
				\coordinate[label=left:$A$] (A) at  (0,0);
				\coordinate[label=below left:$B$]  (B) at (1,-1);
				\coordinate[label=right:$C$] (C) at  (4,0);
				 \coordinate (H) at ($(B)!.5!(C)$);
				\path (1,-2)  node[below]{a)};
				\coordinate[label=left:$A'$] (A1) at ($(H)+(90:4)$); 
				\coordinate[label=left:$B'$] (B1) at  ($(B)-(A)+(A1)$);
				\coordinate[label=right:$C'$] (C1) at  ($(C)-(A)+(A1)$);
				\draw (A1)--(A)--(B)--(C)--(C1)--(A1)--(B1)--(C1) 
					(B)--(B1); 
				\draw[dashed] (A)--(C); 
				\fill (A)  circle (1pt)
					(B)  circle (1pt)
					(C) circle (1pt) 
					 (A1) circle (1pt)
					(B1) circle (1pt) 
					(C1)  circle (1pt);
				\coordinate[label=below  left:$B$] (B2) at (8,-1);
				\coordinate[label=left:$A$] (A2)  at (9,0);
				\coordinate[label=below  right:$C$] (C2) at (12,-1);
				\coordinate[label=right:$D$] (D2) at ($(C2)-(B2)+(A2)$); 
				\coordinate (H1) at  ($(B2)!1/3!(D2)$);
				\coordinate[label=left:$A'$] (A3) at  ($(H1)+(90:4)$);
				\coordinate[label=left:$B'$] (B3) at  ($(B2)-(A2)+(A3)$);
				\coordinate[label=below right:$C'$] (C3) at  ($(C2)-(A2)+(A3)$);
				\coordinate[label=right:$D'$]  (D3) at ($(D2)-(A2)+(A3)$);
				 \path (9,-2) node[below]{b)};
				 \draw (B3)--(B2)--(C2)--(D2)--(D3)--(A3)--(B3)--(C3)--(D3)
					 (C2)--(C3);
				\draw[dashed] (C2)--(A2)--(D2)--(B2) 
					 (A3)--(A2)--(B2);
				\fill (A2)  circle (1pt)
					(B2)  circle (1pt)
					(C2) circle (1pt) 
					(D2) circle (1pt) 
					(A3) circle (1pt) 
					(B3) circle (1pt) 
					 (C3) circle (1pt)
					(D3)  circle (1pt);
				 \coordinate[label=below left:$B$] (B4) at (16,-1);
				\coordinate[label=left:$I$] (A4)  at (17,0);
				 \coordinate[label=below right:$C$] (C4) at (20,-1);
				\coordinate[label=right:$J$] (D4)  at ($(C4)-(B4)+(A4)$);
				\coordinate[label=below:$E$] (E4) at (19,.6); 
				\coordinate  (H2) at ($(B4)!1/3!(D4)$);
				\coordinate[label=left:$I'$] (A5) at ($(H2)+(90:4)$); 
				\coordinate[label=left:$B'$] (B5) at ($(B4)-(A4)+(A5)$); 
				\coordinate[label=above:$C'$] (C5)  at ($(C4)-(A4)+(A5)$);
				\coordinate[label=right:$J'$] (D5)  at ($(D4)-(A4)+(A5)$);
				\coordinate[label=above:$E'$] (E5)  at ($(E4)-(A4)+(A5)$);
				\path  (17,-2) node[below]{c)};
				\draw (B5)--(B4)--(C4)--(D4)--(D5)--(E5)--(A5)--(B5)--(C5)--(D5) 
					 (C4)--(C5);
				\draw[dashed] (B4)--(A4)--(E4)--(D4) 
					(A4)--(A5) 
					(E4)--(E5); 
				\fill (A4)  circle (1pt)
					 (B4) circle (1pt)
					(C4) circle  (1pt)
					(D4) circle (1pt) 
					(E4)  circle (1pt)
					(A5)  circle (1pt)
					(B5) circle (1pt) 
					(C5) circle (1pt) 
					 (D5) circle (1pt)
					 (E5) circle (1pt);
			\end{tikzpicture}
	\end{center}
	 \loigiai{
		\begin{enumerate}
			\item  Hình a) là hình lăng trụ tam giác $ABC.A'B'C'$.\\
			Hình b) là  hình lăng trụ tứ giác $ABCD.A'B'C'D'$.\\
			Hình c) là hình lăng trụ ngũ giác $IBCJE.I'B'C'J'E'$. 
			\item Hình lăng trụ tam giác $ABC.A'B'C'$ trong  hình a) có\\
			Hai mặt đáy là các tam giác  $ABC,A'B'C'$.\\
			Sáu đỉnh $A$, $B$, $C$, $A'$, $B'$,  $C'$.\\
			Ba  mặt bên là các hình bình hành $AA'B'B$, $BB'C'C$, $CC'A'A$.\\
			Ba cạnh bên  $AA'$, $BB'$, $CC'$.
	\end{enumerate}}
\end{ex}

% nguồn: Ok_TuanLuan_Toan11_Lan2
\begin{ex}%[1H4N5-1]%[ma:2a7da0a9]
	Hãy liệt kê các cặp mặt đối diện, các cặp cạnh đối diện và các cặp đỉnh đối diện  của hình hộp $ABCD.A'B'C'D'$.
	\loigiai{ 
		\begin{center}
			\begin{tikzpicture}[line join=round,line cap=round,line width=.6pt,font=\footnotesize,scale=1]
					\coordinate[label=below left:$B$] (B) at  (0,0);
					\coordinate[label=left:$A$] (A)  at (1,.8);
					\coordinate[label=below  right:$C$] (C) at (4,0);
					 \coordinate[label=right:$D$] (D) at ($(C)-(B)+(A)$);
					\coordinate (H) at  ($(B)!1/3!(D)$);
					\coordinate[label=left:$A'$] (A1) at ($(H)+(90:4)$); 
					 \coordinate[label=left:$B'$] (B1) at ($(B)-(A)+(A1)$);
					\coordinate[label=below right:$C'$]  (C1) at ($(C)-(A)+(A1)$);
					\coordinate[label=right:$D'$] (D1) at  ($(D)-(A)+(A1)$);
					\draw (B1)--(B)--(C)--(D)--(D1)--(A1)--(B1)--(C1)--(D1) 
						(C)--(C1); 
					\draw[dashed] (A)--(D) 
						 (A1)--(A)--(B);
					\fill (A)  circle (1pt)
						 (B) circle (1pt)
						(C)  circle (1pt)
						 (D) circle (1pt)
						(A1) circle (1pt) 
						(B1)  circle (1pt)
						 (C1) circle (1pt)
						(D1) circle  (1pt);
				\end{tikzpicture}
		\end{center}
		Trong hình hộp $ABCD.A'B'C'D'$  có\\
		Ba cặp mặt  đối diện là $\left(ABCD\right)$ và $\left(A'B'C'D'\right)$; $\left(ABB'A'\right)$ và $\left(DCC'D'\right)$; $\left(ADD'A'\right)$ và $\left(BCC'B'\right)$.\\
		Sáu cặp cạnh đối diện là $AB$ và $D'C'$; $BC$ và $A'D'$; $CD$ và $B'A'$; $DA$ và $C'B'$; $AA'$ và  $CC'$; $BB'$ và $DD'$.\\
	Bốn cặp đỉnh  đối diện là $A$ và $C'$; $B$ và $D'$; $C$ và $A'$; $D$ và $B'$.}
\end{ex}

% nguồn: Ok_TuanLuan_Toan11_Lan2
\begin{ex}%[1H4H5-1]%[ma:c8246bd9]
	Chứng minh  rằng bốn đường chéo của hình hộp cắt nhau tại trung điểm mỗi đường.
	 \loigiai{
		\begin{center}
			\begin{tikzpicture}[line join=round,line cap=round,line width=.6pt,font=\footnotesize,scale=1]
					\coordinate[label=below  left:$B$] (B) at (0,0);
					\coordinate[label=left:$A$] (A)  at (1,.8);
					\coordinate[label=below right:$C$]  (C) at (4,0);
					\coordinate[label=right:$D$] (D)  at ($(C)-(B)+(A)$);
					\coordinate (H) at  ($(B)!1/3!(D)$);
					 \coordinate[label=left:$A'$] (A1) at ($(H)+(90:4)$);
					 \coordinate[label=left:$B'$] (B1) at ($(B)-(A)+(A1)$);
					\coordinate[label=below right:$C'$] (C1) at  ($(C)-(A)+(A1)$);
					\coordinate[label=right:$D'$] (D1)  at ($(D)-(A)+(A1)$);
					\draw (B1)--(B)--(C)--(D)--(D1)--(A1)--(B1)--(C1)--(D1) 
						 (C)--(C1)--(A1);
					\draw[dashed]  (C)--(A)--(D)
						 (A1)--(A)--(B)
						 (A1)--(C)
						(A)--(C1) 
						(B)--(D1) 
						 (B1)--(D);
					 \fill (A) circle (1pt)
						 (B) circle (1pt)
						 (C) circle (1pt)
						(D) circle  (1pt)
						(A1) circle  (1pt)
						(B1) circle  (1pt)
						 (C1) circle (1pt)
						(D1) circle  (1pt);
				\end{tikzpicture}
		\end{center}
		Cho hình hộp $ABCD.A'B'C'D'$ có các đường  chéo $A'C;AC';BD'$ và $B'D$.\\
		Tứ  giác $ACC'A'$ có $AA'\parallel CC'$ và $AA'=CC'$ (tính chất hình hộp) nên tứ giác $ACC'A'$ là hình bình hành.\\
		Gọi $I$ là giao điểm của hai đường chéo  $AC'$ ' và $A'C$.\\
		Khi đó $I$ là trung điểm của  mỗi đường chéo $AC'$ và $A'C$.\\
		Tương tự, hai tứ giác $A'B'CD$ và $BCD'A'$ cũng là  các hình bình hành nên $I$ là trung điểm của $B'D$ và $BD'$.\\
	Vậy các đường chéo của hình hộp cắt nhau  tại trung điểm mỗi đường.}
\end{ex}
